% Modernised reading — fonds Alexandre Grothendieck, shelfmark 25, pages 121-140
% (only pages 134-135 transcribed).
%
% This is an INTERPRETATION, not a transcription. It re-reads the pages in
% today's notation and vocabulary, states the hypotheses the manuscript leaves
% implicit, and drops the critical apparatus so that the mathematics reads
% straight through. Everything the brackets and underlines carried moves into a
% footnote. It is held to being mathematically correct as it stands; where the
% letter is loose or mistaken, the true statement is what appears here and the
% footnote says what the page has. In any dispute of substance the
% transcription governs: cite batch-07.fr.tex.
%
% Pass: Opus 5.5 (claude-opus-5-5), 2026-10-03 — first pass, unchecked against
% the pages by a human. The model was chosen explicitly for this pass.
%
% Scope. Folder 25 (« Nouveau EGA V », 184 pages) is the typescript that the
% EGA V edition (Blass–Klasa, editions.json id "ega5") covers; per issue #46
% only what that edition does not carry is transcribed: the typed letter to
% Dieudonné of 29.9.1965, pp. 134-135, filed among the typescript of §§ IV
% 19-21 (pp. 121-133, 136-140 of the same batch, not transcribed). This reading
% covers those two pages and nothing else. \pages gives the span of the one
% transcribed batch (121-140), as the skill requires; \pagerange anchors the
% body on 134-135 only. The rest of the folder is NOT restated, and the résumé
% says so.
%
% Spine. One letter, three moves:
%   p. 134      the PLAN of the end of EGA IV as he sees it in September 1965:
%               §§ 20-21 (his former § 20, which Dieudonné is splitting), then
%               §§ 22-27, none of which appeared as such;
%   pp. 134-135 his REFUSAL to sign the « ex-Appendice au par. 18 » alone, and
%               the request that it become a joint paper;
%   p. 135      a STRUCK paragraph giving Dieudonné the limit argument for a
%               statement numbered 10.9.1 in the draft of § 20.
%
% Substantive departure from the page (p. 135, struck paragraph). The letter
% speaks of « l'ensemble des points de Z_λ en lesquels F_λ restreint à la
% fibre est de prof ≥ n » and says it is constructible and even open. Taken
% literally this is false (depth ≥ n fails at every generic point of a fibre
% of positive dimension, and the locus of closed points of the affine line is
% not constructible). The reading states the argument for the loci that §12 of
% EGA IV actually proves open — Serre's (S_n) on the fibres, or « coprofondeur
% ≤ n » — and footnotes the page's wording. The base-change compatibility the
% argument needs (Z_μ = Z_λ ×_{S_λ} S_μ, locus stable under extension of the
% residue field) is supplied and footnoted.
%
% Identifications that are ours, not the folder's: the published titles of
% §§ 20 and 21 of EGA IV (IHÉS 32, 1967); the joint paper Dieudonné–
% Grothendieck, « Critères différentiels de régularité pour les localisés des
% algèbres analytiques », J. Algebra 5 (1967), proposed as what the « ex-
% Appendice au par. 18 » became — plausible by subject (« faire pareil que pour
% les anneaux complets »), not established by anything on these pages. The
% referent of « 10.9.1 » is not identified.
%
% What the letter announces and the folder does not contain: Dieudonné's
% letter « du 24 » and the tables of contents of §§ 16-19 it enclosed; the
% provisional tables of §§ 20-21 asked for; the « vagues notes manuscrites »
% behind the appendix, whose existence the letter itself doubts.
\documentclass[11pt,a4paper]{article}
\input{../preamble/grothendieck}

\folder{25}
\pages{121}{140}
\dating{[à partir de 1967-1987]}
\watermark{Édition de démonstration\\interprétation personnelle de l'œuvre}
\foldertitle{Nouveau EGA V : copies de tapuscrits annotés (s.d.). — lecture modernisée du dossier entier}

\begin{document}

\section*{Résumé}

\begin{resume}
Ce dossier est presque tout entier un tapuscrit : celui d'une suite des
\emph{Éléments de géométrie algébrique}, le grand traité que Grothendieck
écrivait avec Jean Dieudonné, et que l'inventaire appelle « Nouveau EGA V ». Ce
tapuscrit est couvert par une édition existante, celle des prénotes d'EGA V, et
il n'est ni transcrit ni redit ici. Une seule pièce du dossier échappe à cette
édition, et c'est elle que l'on lit : une lettre de deux pages, tapée à la
machine, datée et signée à la main, que Grothendieck adresse à Dieudonné le
29 septembre 1965, et qui s'est retrouvée classée au milieu des feuillets.

Pour comprendre la lettre, il faut savoir comment se fabriquait le traité.
Grothendieck fournissait la matière — des notes, des esquisses, souvent une
phrase dite de vive voix — et Dieudonné rédigeait, numérotait, mettait au
propre, des centaines de pages par an. Le chapitre IV, consacré à l'étude
locale des schémas et des morphismes, paraissait en fascicules successifs, de
1964 à 1967 ; le quatrième et dernier était en rédaction. La lettre est un moment de ce chantier, vu
de l'atelier.

Elle fait trois choses. D'abord, elle remet de l'ordre dans un plan qui
s'échappe : « je commence à me perdre dans le plan », écrit Grothendieck, et il
dresse la liste des paragraphes à venir, de 20 à 27 — systèmes linéaires et
groupe de Picard, grassmanniennes, singularités quadratiques ordinaires,
sections hyperplanes, résultant et discriminant, extensions infinitésimales.
C'est un programme de géométrie projective classique, refait dans le langage
des schémas. Aucun de ces paragraphes n'a paru sous ce numéro : le traité s'est
arrêté au paragraphe 21. Les sujets, eux, ont trouvé place ailleurs, et la
liste est un document sur ce que le traité aurait été.

Ensuite, elle refuse un honneur. Dieudonné a rédigé un appendice et veut le
publier sous le seul nom de Grothendieck. Celui-ci répond que la rédaction ne
doit presque rien à ce qu'il a fourni — « si même je t'en ai jamais passé » ;
il s'était borné à dire qu'il « n'y a qu'à faire pareil que pour les anneaux
complets ». Il ne veut pas non plus que le travail se perde, et propose un
article signé des deux. C'est un trait de la manière de Grothendieck en ces
années : l'idée compte moins à ses yeux que le travail qui la rend utilisable,
et il ne signe pas seul ce qu'un autre a mis au point.

Enfin, un paragraphe qu'il a barré à l'encre donne à Dieudonné un argument
mathématique, et l'accompagne d'un aveu : « C'est vraiment toujours le même
argument qui revient ! » L'argument est celui du \emph{passage à la limite}.
Pour démontrer qu'un objet a une propriété, on l'écrit comme limite d'objets
plus simples — de type fini, par exemple —, on observe que l'ensemble des
points où la propriété vaut est décrit par un nombre fini de conditions, et
l'on conclut qu'elle vaut déjà à un étage fini de l'approximation. C'est
l'analogue algébrique d'un argument de compacité : une suite de fermés
décroissants dont l'intersection est vide l'est déjà à un rang fini. Cet
argument, qui occupe tout le paragraphe 8 du chapitre IV, est l'un des outils
les plus employés du traité, et l'on comprend en le lisant ici pourquoi il
revenait sans cesse sous la plume.

Les noms sous lesquels on cherchera la suite : le chapitre IV des EGA et son
paragraphe 8 (passage à la limite projective, ensembles constructibles), la
condition $(S_n)$ de Serre et la profondeur, le lieu d'ouverture des propriétés
des fibres d'un morphisme plat ; et, pour le plan, les théorèmes de Bertini, les
pinceaux de Lefschetz et les points doubles ordinaires, les grassmanniennes, le
résultant et le discriminant.

\keywords{EGA IV, limit argument, noetherian approximation, constructible set, depth, Serre condition (Sn), flat finitely presented module, openness of fibre loci, linear system, Picard group, Grassmannian, ordinary double point, hyperplane section, resultant, discriminant, infinitesimal extension, criteria of regularity for analytic algebras}
\end{resume}

\pagerange{134}{135}
\section*{Le fil du dossier, et les conventions}

\emph{Ce que contient le dossier, et ce qui est lu ici.} Le dossier 25 compte
184 pages. Elles sont, à cette lettre près, le tapuscrit annoté d'une suite
d'EGA, que couvre l'édition des prénotes d'EGA V ; on ne le redit pas. Seules
les pages 121 à 140 ont fait l'objet d'une transcription, et encore
partiellement : les pages 121 à 133 et 136 à 140 y sont décrites comme le
tapuscrit des paragraphes IV 19 à 21 et la table du « IV 21 », annotés à
l'encre, et laissées à l'édition. La présente lecture porte sur les pages 134
et 135, et sur elles seules\footnote{Les deux feuillets de la lettre portent
au crayon les numéros 37 et 38, de la main du dactylographe ou de la sienne :
la lettre a été rangée dans la liasse des notes, et c'est ce qui explique sa
présence ici.}.

\emph{Une date hors de l'intervalle du dossier.} L'inventaire date le dossier
« à partir de 1967-1987 » ; la lettre porte, de sa main, la date du 29.9.1965.
Elle est donc antérieure à la borne basse, ce qui ne contredit pas l'inventaire
— un dossier constitué à partir de 1967 peut recueillir une pièce plus ancienne
— mais doit être su de qui la cite.

\emph{La lettre en trois mouvements.}
\begin{itemize}
  \item le plan de la fin du chapitre IV (page 134) ;
  \item l'appendice au paragraphe 18 et la proposition d'un article commun
  (pages 134 et 135) ;
  \item un paragraphe barré, qui donne l'argument de limite pour un énoncé du
  paragraphe 20 (page 135).
\end{itemize}

\emph{Conventions.} « Par. » est écrit « § ». Les numéros de paragraphes sont
ceux du chapitre IV d'EGA, que la lettre ne nomme pas autrement que par
« fascicule 4 du Chap IV ». Le tutoiement, et le ton de la lettre — « sans trop
déconner », « et bordel » —, sont conservés tels quels dans les citations ; ils
appartiennent au document.

\pagerange{134}{134}
\section*{I. Le plan de la fin du chapitre IV (page 134)}

Dieudonné a écrit le 24 septembre, en joignant la table des matières des
§§ 16 à 19\footnote{Cette lettre de Dieudonné et la table qu'elle joignait ne
sont pas dans les pages transcrites.}. Grothendieck l'en remercie, demande à
recevoir la table provisoire des §§ 20 et 21, et accepte qu'ils soient joints
au quatrième fascicule du chapitre IV. Puis il pose une question : comment
Dieudonné compte-t-il subdiviser « mon ancien § 20 », et quels seront les
titres des deux morceaux\footnote{La page porte « les », barré à la machine, et
« mon ancien par.20 » ajouté entre les lignes ; la question se termine par un
« 2 » que la transcription lit pour « ? ».} ?

C'est la clef de la numérotation. Dans le plan de Grothendieck, un seul
paragraphe venait après le § 19 ; Dieudonné en fait deux, et ce sont les §§ 20
et 21 publiés dans le quatrième fascicule (1967), « Fonctions méromorphes,
pseudo-morphismes » et « Diviseurs »\footnote{Ces titres sont ceux de la
publication, que la lettre ignore encore : elle écrit « 20.~??? » et
« 21.~??? ». L'identification est la nôtre ; elle s'accorde avec le fait que
les §§ 20 et 21 ferment le fascicule et le chapitre.}. D'où le décalage de
tout ce qui suit, et le besoin de fixer un plan auquel on puisse se référer
sans se tromper de numéro. Le voici, tel qu'il le propose à l'accord de
Dieudonné\footnote{L'orthographe de la page est rectifiée : « Systêmes »,
« Grassmaniennes ».} :

\begin{quote}
20. ? \quad 21. ?\\
22. Systèmes linéaires, compléments sur le groupe de Picard.\\
23. Grassmanniennes.\\
24. Formes lisses, singularités quadratiques ordinaires.\\
25. Sections hyperplanes et bordel.\\
26. Résultant et discriminant.\\
27. Extensions infinitésimales.
\end{quote}

Il ajoute que le § 25 risque d'être fort long, et que Dieudonné voudra sans
doute le couper en deux ; mais que $27 = 3^3$ est « un bien joli nombre ».

Rien de ce programme n'a paru comme paragraphe du chapitre IV : celui-ci
s'achève au § 21, et les EGA avec lui. Le contenu de la liste n'est pas perdu
pour autant, et l'on peut dire où l'on en retrouve l'essentiel\footnote{Ces
renvois sont les nôtres, et indicatifs ; la lettre ne dit rien de ce que
deviendrait le plan.} : les grassmanniennes sont traitées dans la seconde
édition d'EGA I (1971) ; les extensions infinitésimales d'algèbres l'avaient
déjà été dans les compléments du chapitre $0_{\mathrm{IV}}$ (1964) et le seront
plus tard, pour les schémas, par la théorie du complexe cotangent ; les formes
quadratiques, les points doubles ordinaires et les sections hyperplanes en
pinceau forment la matière de SGA 7 (1967--1969), à laquelle répond le
« formes lisses, singularités quadratiques ordinaires » du § 24. Si le reste de ce
dossier, la suite d'EGA qu'il contient, reprend une part de ce programme, c'est
à l'édition qui le couvre de le dire ; on ne le vérifie pas ici, puisque ce
reste n'est pas transcrit.

\pagerange{134}{135}
\section*{II. L'appendice au § 18, et un article à deux (pages 134 et 135)}

Dieudonné a rédigé un texte, l'« ex-Appendice au § 18 », et veut le publier
sous le nom de Grothendieck. Celui-ci refuse : la rédaction n'a « à peu près
plus rien de commun » avec les notes manuscrites qu'il lui aurait passées
— « si même je t'en ai jamais passé » — et lui-même s'est borné à dire « il n'y
a qu'à faire pareil que pour les anneaux complets\footnote{« manuscriptes »
sur la page.} ». Mais il serait dommage que ce travail de mise au point soit
perdu pour ses utilisateurs éventuels, « il finit toujours par s'en trouver » ;
il demande donc à Dieudonné de reconsidérer la question d'en faire un « joint
paper\footnote{La page porte « papar », que la transcription lit pour
« paper ».} ».

Ce que fut cet appendice, la lettre ne le dit pas. Le paragraphe 18 du
chapitre IV traite des morphismes étales et des anneaux henséliens ; les
anneaux complets ont reçu, dans les compléments du chapitre $0_{\mathrm{IV}}$,
des critères différentiels de régularité. « Faire pareil » pour une autre
classe d'anneaux locaux henséliens désigne naturellement les algèbres
analytiques, et un article signé des deux auteurs a paru sur ce sujet en
1967 : « Critères différentiels de régularité pour les localisés des algèbres
analytiques » (\emph{Journal of Algebra} 5)\footnote{L'identification est
la nôtre, et elle n'est qu'une conjecture plausible : elle repose sur la
coïncidence du sujet, de la date et de la double signature, et sur rien qui
soit écrit dans le dossier.}. Si elle est juste, la lettre est l'origine de cet
article, et la proposition qu'elle fait a été acceptée.

\pagerange{135}{135}
\section*{III. Le paragraphe barré : l'argument de limite (page 135)}

Avant la formule de politesse, un paragraphe donne à Dieudonné un indice pour
un énoncé du futur § 20, que la lettre numérote 10.9.1\footnote{Le paragraphe
entier est barré d'un trait oblique à l'encre, après coup : Grothendieck l'a
retiré de la lettre, sans qu'on sache s'il l'a fait avant ou après envoi, ni
s'il a transmis l'indication autrement. Le numéro « 10.9.1 » renvoie
vraisemblablement à la rédaction en cours du § 20 ; nous ne l'avons pas
identifié.}. Sa forme est télégraphique ; voici l'argument qu'il désigne, sous
les hypothèses qu'il suppose sans les dire.

\emph{La situation.} On a un système projectif filtrant $(S_\lambda)$ de
schémas quasi compacts et quasi séparés, à morphismes de transition affines,
de limite $S$ ; un schéma $Z_\lambda$ de présentation finie sur l'un d'eux,
$S_\lambda$, et un $\mathcal{O}_{Z_\lambda}$-module $F_\lambda$ de présentation
finie et plat sur $S_\lambda$ ; on pose, pour $\mu \geq \lambda$,
$Z_\mu = Z_\lambda \times_{S_\lambda} S_\mu$, $Z = Z_\lambda \times_{S_\lambda}
S$, et l'on note $F_\mu$, $F$ les images réciproques de
$F_\lambda$\footnote{La lettre parle de « $Z_\lambda$ », de « $F_\lambda$ »
et de « son image inverse dans $Z$ », sans plus. Les hypothèses de finitude et
le fait que $Z_\mu$ s'obtienne par changement de base sont ajoutés : sans eux,
ni la constructibilité ni la compatibilité invoquées plus bas ne sont
garanties. Ce sont les hypothèses du § 8 du chapitre IV, et la lettre y fait
allusion en rappelant « les hypothèses de platitude et de présentation finie
qu'on a faites ».}.

\emph{Le lieu.} Fixons un entier $n$ et appelons $E_\lambda$ l'ensemble des
points $z$ de $Z_\lambda$ où la restriction de $F_\lambda$ à la fibre de $z$
satisfait la condition $(S_n)$ de Serre : en notant $s$ l'image de $z$
dans $S_\lambda$ et $M = (F_\lambda \otimes k(s))_z$, module sur l'anneau
local de $z$ dans sa fibre, on demande
\[ \operatorname{prof} M_{\mathfrak{p}} \;\geq\;
\inf\bigl(n,\ \dim M_{\mathfrak{p}}\bigr) \]
pour tout idéal premier $\mathfrak{p}$ de cet anneau local — ou bien, variante, l'ensemble des
points où la coprofondeur de la fibre est au plus $n$\footnote{La page dit :
l'ensemble des points en lesquels $F_\lambda$ restreint à la fibre « est de
prof $\geq n$ » (le signe $\geq$ est d'une lecture incertaine). Prise à la
lettre, la condition « profondeur au moins $n$ » ne définit pas un ensemble
constructible : sur la droite affine au-dessus d'un corps, avec
$F = \mathcal{O}$ et $n = 1$, elle est vérifiée exactement aux points fermés,
et l'ensemble des points fermés n'est pas constructible, le point générique
n'étant pas, seul, un ensemble constructible. Les lieux que le § 12 du
chapitre IV démontre ouverts, sous les hypothèses de platitude et de
présentation finie, sont ceux de la condition $(S_n)$ et de la coprofondeur
bornée sur les fibres ; c'est d'eux que l'argument a besoin, et c'est
probablement ce que « prof $\geq n$ » abrège.}. Deux faits sont acquis :
\begin{itemize}
  \item $E_\lambda$ est \emph{constructible}, et même \emph{ouvert} — c'est ce
  que le § 12 a établi, sous les hypothèses de platitude et de présentation
  finie (la lettre : « on a même dû prouver au par.~12 qu'il est ouvert ») ;
  \item la formation de ce lieu commute au changement de base : l'image
  réciproque de $E_\lambda$ dans $Z_\mu$ est $E_\mu$, et son image réciproque
  dans $Z$ est le lieu correspondant pour $F$. Cela tient à ce que les fibres
  de $Z_\mu$ se déduisent de celles de $Z_\lambda$ par extension du corps de
  base, et que la condition $(S_n)$ d'un module cohérent sur un schéma
  localement de type fini sur un corps est insensible à une telle
  extension\footnote{Ce second point n'est pas dans la lettre ; il est
  nécessaire pour que la conclusion porte sur $F_\mu$ et non seulement sur
  l'image de $Z_\mu$ dans $Z_\lambda$.}.
\end{itemize}

\emph{L'argument.} Supposons qu'au niveau limite le lieu soit tout $Z$ ;
c'est l'hypothèse de l'énoncé. Alors le fermé constructible
$Z_\lambda \smallsetminus E_\lambda$ a une image réciproque vide dans $Z$. Or,
d'après le § 8 du chapitre IV, une partie constructible de $Z_\lambda$ dont
l'image réciproque dans la limite est vide a déjà une image réciproque vide
dans $Z_\mu$ pour un $\mu \geq \lambda$ assez grand. C'est l'analogue d'un
argument de compacité : les images réciproques forment une famille filtrante
décroissante de constructibles d'intersection vide, et la topologie
constructible est quasi compacte. Pour ce $\mu$, on a donc $E_\mu = Z_\mu$ :
la propriété, vraie à la limite, est vraie « un peu plus loin que
$\lambda$ ».

Le paragraphe s'achève sur la remarque qui en fait le prix : « C'est vraiment
toujours le même argument qui revient ! » Le § 8 du chapitre IV a été écrit
pour cela — pour qu'un énoncé démontré sur des schémas noethériens, ou sur des
schémas de type fini, s'étende à des situations de présentation finie par
approximation, une propriété à la fois. La lettre en montre l'usage le plus
simple : constructibilité du lieu, compatibilité au changement de base, et
quasi-compacité, en trois lignes.

\end{document}
